\section{Conclusion}
\label{sec:conclusion}

Three in four Albanian fifteen-year-olds were below basic mathematics proficiency in
2022, and the number with which we began now has an explanation. The collapse was not a
uniformly weaker cohort but a change in the structure of who is at risk and why: a
domain classifier separates the 2018 and 2022 cohorts almost perfectly from background
data alone, the variable that moved most was the support students receive in school
rather than the resources they have at home, and the decline sits on the flattest
socioeconomic gradient of the nine countries we compare. Read against Albania's
history, a destructive 2019 earthquake, an extended pandemic closure, a thin
remote-learning infrastructure, and the steady emigration of teachers and working-age
adults, the statistical shape and the economic record tell the same story: a
supply-side disruption to schooling, not a demand-side collapse in family means.

For anyone who would predict this risk, the practical lesson is where the ceiling comes
from. Individual background answers plateau at about 0.73 AUC and no amount of tuning
moves them; the single missing ingredient is the socioeconomic composition of a
student's school, which lifts every model to about 0.78, after which a multilevel model
and a direct school-questionnaire linkage independently confirm that the ceiling is now
a property of what background questionnaires can measure rather than of the model. A
plateau, in other words, can be a feature ceiling or a data ceiling, and telling them
apart takes convergent evidence rather than one more round of optimisation.

For anyone who would act on it, the lesson is that a screener trained on a single
cross-section reports associations, not levers. The same analysis that ranks math
anxiety highly also shows why reducing a model input is not an intervention estimate,
and the fairness audit shows a naive threshold over-flagging the poorest students. The
tool we release is therefore built to widen support, never to ration it. The policy
direction that follows is not targeted transfers to the poor, whom a flat gradient says
are not the whole problem, but the rebuilding of schooling capacity for everyone:
retaining teachers, reconstructing schools, and funding catch-up instruction.

Two questions that earlier framings left open we were able to take a first pass at here,
and both answers point the same way. Richer data \emph{can} move the 0.78 ceiling,
PISA's within-test process data reaches $\sim$0.86, but only through a signal endogenous
to the score itself, so the deployable ceiling holds and the genuinely exogenous process
records (attendance, formative assessment, longitudinal files) that would capture the
disruption directly remain the route worth building. And the one causal question the
observational data admit, whether the 2019 earthquake localised the collapse, returns a
well-identified null: a difference-in-differences with a failing placebo shows the loss
was national, not concentrated at the epicentre. What remains is a design that can
license a \emph{positive} policy lever rather than rule a single-region story out. Albania's 2022 result is the clearest natural signal yet of what
a compounded shock does to a fragile education system; understanding it well is the
difference between planning for recovery and mistaking a broken process for a weak
generation of students.

\paragraph{Reproducibility.}
Code, configuration, the unit-tested pipeline, the calibrated bilingual screener, and
this manuscript's figures are in the project repository. Raw PISA microdata are public
from the OECD but not redistributed here.
